\documentclass{report}
\usepackage{luatexja} % LuaTeXで日本語を使うためのパッケージ
\usepackage{luatexja-fontspec} % LuaTeX用の日本語フォント設定

% --- 数学関連 ---
\usepackage{amsmath, amssymb, amsfonts, mathtools, bm, amsthm} % 基本的な数学パッケージ
\usepackage{type1cm, upgreek} % 数式フォントとギリシャ文字
\usepackage{physics, mhchem} % 物理や化学の記号や式の表記を簡単にする

% --- 表関連 ---
\usepackage{multirow, longtable, tabularx, array, colortbl, dcolumn, diagbox} % 表のレイアウトを柔軟にする
\usepackage{tablefootnote, truthtable} % 表中に注釈を追加、真理値表
\usepackage{tabularray} % 高度な表組みレイアウト

% --- グラフィック関連 ---
\usepackage{tikz, graphicx} % 図の描画と画像の挿入
\usepackage{background} % ウォーターマークの設定
\usepackage{caption, subcaption} % 図や表のキャプション設定
\usepackage{float, here} % 図や表の位置指定

% --- レイアウトとページ設定 ---
\usepackage{fancyhdr} % ページヘッダー、フッター、余白の設定
\usepackage[top = 20truemm, bottom = 20truemm, left = 20truemm, right = 20truemm]{geometry}
\usepackage{fancybox, ascmac} % ボックスのデザイン

% --- 色とスタイル ---
\usepackage{xcolor, color, colortbl, tcolorbox} % 色とカラーボックス
\usepackage{listings, jvlisting} % コードの色付けとフォーマット

% --- 参考文献関連 ---
\usepackage{biblatex, usebib} % 参考文献の管理と挿入
\usepackage{url, hyperref} % URLとリンクの設定

% --- その他の便利なパッケージ ---
\usepackage{footmisc} % 脚注のカスタマイズ
\usepackage{multicol} % 複数段組
\usepackage{comment} % コメントアウトの拡張
\usepackage{siunitx} % 単位の表記
\usepackage{docmute}
% \usepackage{appendix}
% --- tcolorboxとtikzの設定 ---
\tcbuselibrary{theorems, breakable} % 定理のボックスと改ページ設定
\usetikzlibrary{decorations.markings, arrows.meta, calc} % tikzの装飾や矢印の設定

% --- 定理スタイルと数式設定 ---
\theoremstyle{definition} % 定義スタイル
\numberwithin{equation}{section} % 式番号をサブセクション単位でリセット

% --- hyperrefの設定 ---
\hypersetup{
  setpagesize = false,
  bookmarks = true,
  bookmarksdepth = tocdepth,
  bookmarksnumbered = true,
  colorlinks = false,
  pdftitle = {}, % PDFタイトル
  pdfsubject = {}, % PDFサブジェクト
  pdfauthor = {}, % PDF作者
  pdfkeywords = {} % PDFキーワード
}

% --- siunitxの設定 ---
\sisetup{
  table-format = 1.5, % 小数点以下の桁数
  table-number-alignment = center, % 数値の中央揃え
}

% --- 透かし画像の設定 --- 
\backgroundsetup{
  scale=0.5,                       % 画像のスケール
  % color=black,                   % 画像の色（透かし用に半透明が推奨）
  opacity=0.2,                   % 透かしの透明度（0が完全透明、1が完全不透明）
  angle=0,                       % 画像の角度
  position = current page.south east,  % ページの右下
  hshift=-6cm, % 右方向へのシフト（負の値で内側に移動）
  vshift=5cm,  % 上方向へのシフト（正の値で内側に移動）
  contents={\includegraphics{./fig/appilogo-circular-full.png}} % 画像のパス
}

% --- その他の設定 ---
\allowdisplaybreaks % 数式の途中改ページ許可
\newcolumntype{t}{!{\vrule width 0.1pt}} % 新しいカラムタイプ
\newcolumntype{b}{!{\vrule width 1.5pt}} % 太いカラム
\UseTblrLibrary{amsmath, booktabs, counter, diagbox, functional, hook, html, nameref, siunitx, varwidth, zref} % tabularrayのライブラリ
\setlength{\columnseprule}{0.4pt} % カラム区切り線の太さ
\captionsetup[figure]{font = bf} % 図のキャプションの太字設定
\captionsetup[table]{font = bf} % 表のキャプションの太字設定
\captionsetup[lstlisting]{font = bf} % コードのキャプションの太字設定
\captionsetup[subfigure]{font = bf, labelformat = simple} % サブ図のキャプション設定
\setcounter{secnumdepth}{4} % セクションの深さ設定
\newcolumntype{d}{D{.}{.}{5}} % 数値のカラム
\newcolumntype{M}[1]{>{\centering\arraybackslash}m{#1}} % センター揃えのカラム
\DeclareMathOperator{\diag}{diag}
\everymath{\displaystyle} % 数式のスタイル
\newcommand{\inner}[2]{\left\langle #1, #2 \right\rangle}
\renewcommand{\figurename}{図}
\renewcommand{\i}{\mathrm{i}} % 複素数単位i
\renewcommand{\laplacian}{\grad^2} % ラプラシアンの記号
\renewcommand{\thesubfigure}{(\alph{subfigure})} % サブ図の番号形式
\newcommand{\m}[3]{\multicolumn{#1}{#2}{#3}} % マルチカラムのショートカット
\renewcommand{\r}[1]{\mathrm{#1}} % mathrmのショートカット
\newcommand{\e}{\mathrm{e}} % 自然対数の底e
\newcommand{\Ef}{E_{\mathrm{F}}} % フェルミエネルギー
\renewcommand{\c}{\si{\degreeCelsius}} % 摂氏記号
\renewcommand{\d}{\r{d}} % d記号
\renewcommand{\t}[1]{\texttt{#1}} % タイプライタフォント
\newcommand{\kb}{k_{\mathrm{B}}} % ボルツマン定数
% \renewcommand{\phi}{\varphi} % ϕをφに変更
\renewcommand{\epsilon}{\varepsilon}
\newcommand{\fullref}[1]{\textbf{\ref{#1} \nameref{#1}}}
\newcommand{\reff}[1]{\textbf{図\ref{#1}}} % 図参照のショートカット
\newcommand{\reft}[1]{\textbf{表\ref{#1}}} % 表参照のショートカット
\newcommand{\refe}[1]{\textbf{式\eqref{#1}}} % 式参照のショートカット
\newcommand{\refp}[1]{\textbf{コード\ref{#1}}} % コード参照のショートカット
\renewcommand{\lstlistingname}{コード} % コードリストの名前
\renewcommand{\theequation}{\thesection.\arabic{equation}} % 式番号の形式
\renewcommand{\footrulewidth}{0.4pt} % フッターの線
\newcommand{\mar}[1]{\textcircled{\scriptsize #1}} % 丸囲み文字
\newcommand{\combination}[2]{{}_{#1} \mathrm{C}_{#2}} % 組み合わせ
\newcommand{\thline}{\noalign{\hrule height 0.1pt}} % 細い横線
\newcommand{\bhline}{\noalign{\hrule height 1.5pt}} % 太い横線

% --- カスタム色定義 ---
\definecolor{burgundy}{rgb}{0.5, 0.0, 0.13} % バーガンディ色
\definecolor{charcoal}{rgb}{0.21, 0.27, 0.31} % チャコール色
\definecolor{forest}{rgb}{0.0, 0.35, 0} % 森の緑色

% --- カスタム定理環境の定義 ---
\newtcbtheorem[number within = chapter]{myexc}{練習問題}{
  fonttitle = \gtfamily\sffamily\bfseries\upshape,
  colframe = forest,
  colback = forest!2!white,
  rightrule = 1pt,
  leftrule = 1pt,
  bottomrule = 2pt,
  colbacktitle = forest,
  theorem style = standard,
  breakable,
  arc = 0pt,
}{exc-ref}
\newtcbtheorem[number within = chapter]{myprop}{命題}{
  fonttitle = \gtfamily\sffamily\bfseries\upshape,
  colframe = blue!50!black,
  colback = blue!50!black!2!white,
  rightrule = 1pt,
  leftrule = 1pt,
  bottomrule = 2pt,
  colbacktitle = blue!50!black,
  theorem style = standard,
  breakable,
  arc = 0pt
}{proposition-ref}
\newtcbtheorem[number within = chapter]{myrem}{注意}{
  fonttitle = \gtfamily\sffamily\bfseries\upshape,
  colframe = yellow!20!black,
  colback = yellow!50,
  rightrule = 1pt,
  leftrule = 1pt,
  bottomrule = 2pt,
  colbacktitle = yellow!20!black,
  theorem style = standard,
  breakable,
  arc = 0pt
}{remark-ref}
\newtcbtheorem[number within = chapter]{myex}{例題}{
  fonttitle = \gtfamily\sffamily\bfseries\upshape,
  colframe = black,
  colback = white,
  rightrule = 1pt,
  leftrule = 1pt,
  bottomrule = 2pt,
  colbacktitle = black,
  theorem style = standard,
  breakable,
  arc = 0pt
}{example-ref}
\newtcbtheorem[number within = chapter]{exc}{Requirement}{myexc}{exc-ref}
\newcommand{\rqref}[1]{{\bfseries\sffamily 練習問題 \ref{exc-ref:#1}}}
\newtcbtheorem[number within = chapter]{definition}{Definition}{mydef}{definition-ref}
\newcommand{\dfref}[1]{{\bfseries\sffamily 定義 \ref{definition-ref:#1}}}
\newtcbtheorem[number within = chapter]{prop}{命題}{myprop}{proposition-ref}
\newcommand{\prref}[1]{{\bfseries\sffamily 命題 \ref{proposition-ref:#1}}}
\newtcbtheorem[number within = chapter]{rem}{注意}{myrem}{remark-ref}
\newcommand{\rmref}[1]{{\bfseries\sffamily 注意 \ref{remark-ref:#1}}}
\newtcbtheorem[number within = chapter]{ex}{例題}{myex}{example-ref}
\newcommand{\exref}[1]{{\bfseries\sffamily 例題 \ref{example-ref:#1}}}
% --- 再定義コマンド ---
% \mathtoolsset{showonlyrefs=true} % 必要な式番号のみ表示
\pagestyle{fancy} % ヘッダー・フッターのスタイル設定
\chead{応用量子物性講義ノート} % 中央ヘッダー
% \rhead{}
\fancyhead[R]{\rightmark}
\renewcommand{\sectionmark}[1]{\markright{\thesection\ #1}}
\cfoot{\thepage} % 中央フッターにページ番号
\lhead{}
\rfoot{Yuto Masuda and Haruki Aoki} % 右フッターに名前
\setcounter{tocdepth}{4} % 目次の深さ
\makeatletter
\@addtoreset{equation}{section} % サブセクションごとに式番号をリセット
\makeatother

% --- メタ情報 ---
\title{応用量子物性講義ノート}
\date{更新日\today}
\author{Yuto Masuda and Haruki Aoki}

\begin{document}
  1次摂動のエネルギー補正を求めるときに縮退が無いという条件である\refe{no-degeneracy}を用いて，
  \begin{align}
    \qty(E_n^{(0)} - \hat{H}^{(0)})\ket{n^{(1)}} + E_n^{(1)}\ket{n^{(0)}} &= \hat{V}\ket{n^{(0)}} \label{1st-order-perturbation-re-before} \\
    \ket{n^{(1)}} &= \sum_{m\neq n}\ket{m^{(0)}}\frac{\mel**{m^{(0)}}{\hat{V}}{n^{(0)}}}{E_n^{(0)} - E_m^{(0)}}\label{1st-order-perturbation-re-after}
  \end{align}
  と計算していた．
  では，エネルギー縮退が存在するとき，どのように計算すればよいだろうか．
  $\hat{H}^{(0)}$の固有値$E_n^{(0)}$に，異なる2つの固有ベクトル$\ket{n_a^{(0)}}$，$\ket{n_b^{(0)}}$が属すると状況を考える．
  すなわち，
  \begin{align}
    \hat{H}^{(0)}\ket{n_a^{(0)}} &= E_n^{(0)}\ket{n_a^{(0)}} \\
    \hat{H}^{(0)}\ket{n_b^{(0)}} &= E_n^{(0)}\ket{n_b^{(0)}}
  \end{align}
  が成立するときである．
  ただし，Hermite演算子の固有ベクトルは規格直交化できるので，$\braket{n_i^{(0)}}{n_j^{(0)}} = \delta_{ij}$とする．
  $\ket{n_a^{(0)}}$と$\ket{n_b^{(0)}}$は同じ固有値をもつため，これらの線形結合，
  \begin{align}
    \ket{n^{(0)}} = \alpha\ket{n_a^{(0)}} + \beta\ket{n_b^{(0)}}\label{degeneracy-n0-def}
  \end{align}
  も，固有値$E_n^{(0)}$に属する固有ベクトルである．
  \par
  まず，\refe{1st-order-perturbation-re-before}の両辺に左から$\bra{n_a^{(0)}}$を作用すると，
  \begin{align}
    \mel**{n_a^{(0)}}{\qty(E_n^{(0)} - \hat{H}^{(0)})}{n^{(1)}} + \mel**{n_a^{(0)}}{E_n^{(1)}}{n^{(0)}} &= \mel**{n_a^{(0)}}{\hat{V}}{n^{(0)}} \\ 
    \Leftrightarrow E_n^{(0)}\braket{n_a^{(0)}}{n^{(1)}} - E_n^{(0)}\braket{n_a^{(0)}}{n^{(1)}} + E_n^{(1)}\braket{n_a^{(0)}}{n^{(0)}} &= \mel**{n_a^{(0)}}{\hat{V}}{n^{(0)}} \\ 
    \Leftrightarrow \alpha E_n^{(1)} &= \mel**{n_a^{(0)}}{\hat{V}}{n^{(0)}}\label{degeneracy-ket-alpha}
  \end{align}
  となる．
  $V_{ij} \coloneqq \mel**{n^{(0)}_i}{\hat{V}}{n^{(0)}_j}$とすれば\refe{degeneracy-ket-alpha}は，
  \begin{align}
    \alpha E_n^{(1)} = \alpha V_{aa} + \beta V_{ab}
  \end{align}
  と書ける．同様に\refe{1st-order-perturbation-re-before}の両辺に左から$\bra{n_b^{(0)}}$を作用させたときも考えれば，
  \begin{align}
    \begin{dcases}
      \alpha V_{aa} + \beta V_{ab} = \alpha E_n^{(1)}\\
      \alpha V_{ba} + \beta V_{bb} = \beta E_n^{(1)}
    \end{dcases}
  \end{align}
  を得る．これは行列を用いて，
  \begin{align}
    \mqty(
      V_{aa} - E_n^{(1)} & V_{ab} \\
      V_{ba} & V_{bb} - E_n^{(1)})
    \mqty(\alpha\\ \beta)
    =
    \mqty(0 \\ 0)\label{secular-eq}
  \end{align}
  のように書き直される．
  行列の部分がHermite行列になっているので，\refe{secular-eq}は永年方程式である．
  永年方程式が非自明な解を持つ条件は，
  \begin{align}
    \mqty|
      V_{aa} - E_n^{(1)} &  V_{ab}\\
      V_{ba} & V_{bb} - E_n^{(1)}|
    = 0
  \end{align}
  である．よって1次の摂動エネルギーとして，
  \begin{itembox}[l]{縮退がある場合の摂動論による1次エネルギー補正}
    \begin{align}
      E_n^{(1)} = \frac{1}{2}\qty[(V_{aa} + V_{bb})\pm\sqrt{(V_{aa} - V_{bb})^2 + 4\abs{V_{ab}}^2}]\label{1st-order-energy-with-degeneracy}
    \end{align}
  \end{itembox}
  を得る．
  \refe{1st-order-energy-with-degeneracy}をみると，定常状態では1種類であったエネルギーが，
  縮退が解けて2つに分かれている．
  \begin{myex}{}{}
    $ V_{aa} = V_{bb} = 0$，$ V_{ab} = V_{ba} = V$の場合
    \refe{secular-eq}は，
    \begin{align}
      \mqty(
        -E_n^{(1)} &  V\\
        V & -E_n^{(1)})
      \mqty(\alpha\\\beta)
      = \mqty(0\\ 0)
    \end{align}
    となる．よって1次エネルギー補正は，
    \begin{align}
      E_n^{(1)}=
      \begin{dcases}
        + V & \alpha = 1, \beta = 1\\
        - V & \alpha = 1, \beta = -1
      \end{dcases}
    \end{align}
    と求まる．
    これは摂動を加える前に縮退していた2つの状態$\ket{n^{(0)}} = \ket{n_a^{(0)}}\pm\ket{n_b^{(0)}}$の縮退が解け，
    エネルギー$E_n^{(0)} +  V$をもつ状態$\ket{n_a^{(0)}} + \ket{n_b^{(0)}}$とエネルギー$E_n^{(0)} -  V$をもつ状態$\ket{n_a^{(0)}} - \ket{n_b^{(0)}}$に
    分かれたことを意味している．
  \end{myex}
\end{document}
